% Generated draft: requires mathematical and source review.
\documentclass[11pt]{article}
\usepackage[margin=25mm]{geometry}
\usepackage{fontspec}
\setmainfont{Latin Modern Roman}
\usepackage{amsmath,amssymb,mathtools,unicode-math}
\setmathfont{Latin Modern Math}
\usepackage{enumitem,needspace,xurl,hyperref}
\setlength{\parindent}{0pt}
\setlength{\parskip}{5pt}
\setlength{\emergencystretch}{3em}
\newcommand{\N}{\mathbb{N}}
\newcommand{\Z}{\mathbb{Z}}
\newcommand{\Q}{\mathbb{Q}}
\newcommand{\R}{\mathbb{R}}
\newcommand{\C}{\mathbb{C}}
\newcommand{\F}{\mathbb{F}}
\newcommand{\A}{\mathbb{A}}
\newcommand{\PP}{\mathbb{P}}
\newcommand{\E}{\mathbb{E}}
\newcommand{\eps}{\varepsilon}
\newcommand{\dd}{\,\mathrm{d}}
\begin{document}
\section*{UnsolvedMath ID: 701291}
\textbf{Source problem code:} catalogue entry
\subsection{Definitions and mathematical statement}
Let \(F_2=\langle x,y\rangle\) be the free group of rank two. For a group \(G\), define
\[
\mathcal K_G=\{N\trianglelefteq F_2:F_2/N\cong G\}.
\]
The automorphism group \(\operatorname{Aut}(F_2)\) acts on \(\mathcal K_G\) by \(\alpha\cdot N=\alpha(N)\). A group is two-generated if it is generated by two elements, and it is simple if it is nontrivial and has no normal subgroups other than the identity subgroup and itself.

The infinite version of the question is: does there exist an infinite two-generated simple group \(G\) such that
\[
\forall N,M\in\mathcal K_G\quad\exists\alpha\in\operatorname{Aut}(F_2)\quad\alpha(N)=M?
\]
Two-generation ensures that \(\mathcal K_G\ne\varnothing\), so this transitivity condition is not vacuous.

The source writes ``(infinite)'' before ``2-generator simple group.'' The displayed question explicitly retains infinity as a requirement; if the parentheses instead indicate an optional restriction, the broader question is the same existence question with ``infinite'' omitted. The supplied record does not resolve that distinction. Its comparison with Problem 6.45 is retained as a source qualification, but the statement of that problem is not supplied.
\subsection{Short English statement}
Does an infinite simple group generated by two elements exist for which every two normal subgroups of the rank-two free group giving that quotient are carried to each other by an automorphism of the free group? The source puts ``infinite'' in parentheses, leaving its intended necessity unclear.
\subsection{Sources}
\begin{itemize}
\item Supplied local UnsolvedMath registry record: permanent ID \(701291\), exact source problem code \texttt{catalogue entry}, title ``Kourovka Notebook Problem 16.100.'' The supplied status is \texttt{open}; this is a registry label, not an independently verified mathematical assessment.
\item Attributed in the supplied record to J. Wiegold, proposed in 2006, in \emph{The Kourovka Notebook (Unsolved Problems in Group Theory)}, 16th issue (2006), Problem 16.100. The record also identifies the 21st issue edition (2026), arXiv:1401.0300v48, and names E. I. Khukhro and V. D. Mazurov as editors of the collection. These bibliographic details are supplied attribution and have not been independently verified.
\item The source statement includes ``Cf. 6.45.'' No statement or further details of that comparison problem are supplied.
\end{itemize}
\end{document}
